% Setup of the supporting analyses: in the main text for arXiv, in the online supplement for the journal.
\paragraph{Intervals}
Each model's interval forecasts come from its own errors in the earlier rolling
windows of the same data set, a rolling form of conformal
calibration~\cite{vovk2005}. With $e_{j,h} = (y_{j,t+h} - \hat y_{j,t+h})/\sigma_j$
the scaled errors of the earlier windows at horizon $h$ ($\sigma_j$ the series'
one-step scale), the interval at level $1-\alpha$ is
$\big[\hat y_{j,t+h} + \sigma_j q_{\alpha/2,h},\ \hat y_{j,t+h} + \sigma_j
q_{1-\alpha/2,h}\big]$, with $q_{\cdot,h}$ the empirical quantiles of the pooled
$e_{\cdot,h}$. Where a level has fewer than 50
earlier errors for each horizon, all horizons are pooled into one set of
quantiles, and interval ends are rounded outwards to whole units except on
Favorita; this fix was added on 16 October after the first intervals covered
far less than planned at the totals, and is used throughout. The first
rolling window has no earlier window and is not scored. We report the scaled
pinball loss averaged over the quantiles, weighted as in the M5 uncertainty
competition~\cite{m5uncertainty}, and coverage at 50\%, 80\% and 95\%. Models
that give quantiles of their own (Chronos-Bolt, TFT) are also scored on them,
over every window.

\paragraph{Decisions}
Three simulations turn forecasts into costs. \emph{Store orders}: for each
item-store-period, order the model's interval quantile at a critical ratio $r$
of 0.5, 0.8 or 0.95; each unit over costs 1 and each unit short $r/(1-r)$.
\emph{Category plan}: for category-by-store nodes $n$ with sales $y_{n,t}$,
forecast $\hat y_{n,t}$ and average price $\bar\pi_n$ before the window,
\begin{equation}
C = \frac{\sum_n \bar\pi_n \sum_{t} |y_{n,t} - \hat y_{n,t}|}{\sum_n \bar\pi_n \sum_t y_{n,t}},
\end{equation}
the absolute error in dollars over sales in dollars. \emph{Promotion decisions}: in the generated
store worlds of \ifjournal an audit of these models\else the companion paper\fi~\cite{paper2}, where each promotion's true
lift is known, promote when the predicted lift times margin exceeds a fixed
cost; regret is the value lost against deciding with the true lift. The first
two are reported against seasonal naive. All are simulations under stated
costs, not business results.

\paragraph{Reconciliation}
Base forecasts at every node come from AutoETS fitted on the history before each
window (seasonal at the aggregate nodes, non-seasonal item by store). They are
reconciled bottom-up, top-down by historical proportions, and by the MinT
family~\cite{mint}. With $S$ the summing matrix (nodes by item-store series) and
$\hat{\mathbf y}$ the base forecasts at every node, the reconciled item-store
forecasts are
\begin{equation}
\tilde{\mathbf b} = (S^\top W^{-1} S)^{-1} S^\top W^{-1}\hat{\mathbf y},
\end{equation}
with $W = I$ (ordinary least squares), $W = \mathrm{diag}(S\mathbf 1)$ (structural
weights), $W$ the diagonal of in-sample error variances, or $W = \lambda
\mathrm{diag}(\hat\Sigma) + (1-\lambda)\hat\Sigma$ with the Sch\"afer--Strimmer
intensity $\lambda$ (MinT-shrink, where the hierarchy has at most 6,000 nodes:
Frat); nodes with no
in-sample error are given the smallest variance of the others. Each model's own
item-by-store forecasts are also combined with ETS above them by weighted least
squares.

\paragraph{Ranking}
Accuracy says how close a forecast is, not whether it puts items in the right
order. For every window we report the Spearman correlation, across item-store
series that sold in the window before, between the forecast change (horizon
total over the previous horizon's total) and the actual change; the share of the
tenth of nodes forecast to grow most that are in the actual top tenth (0.10 by
chance); the Spearman correlation across promoted periods between forecast and
actual units over the series' unpromoted level; and bias, forecast over actual.
These were written down before they ran and have no pass rule.
